\section{Torque history and the remaining analytic obstruction}\label{sec:torque}

An analytic restriction on trajectories released from the tied brake curve
could extend the interval results. Integrating the pair torques gives
necessary cancellation conditions for a second brake and a scalar equation
for the accumulated torque ratio. The proposed global sign consequence
remains unproved.

\subsection{Integral constraints}

Let $d_{ij}=q_j-q_i$, $r_{ij}=|d_{ij}|$, and
$\Delta_2=(q_2-q_1)\times(q_3-q_1)$ be twice the signed area. Define the
specific pair angular momenta $\ell_{ij}=d_{ij}\times\dot d_{ij}$, using
the cyclic pairs $12,23,31$. Direct differentiation gives
\begin{equation}\label{eq:torques}
\begin{aligned}
 \dot\ell_{12}&=m_3\Delta_2(r_{23}^{-3}-r_{31}^{-3}),\\
 \dot\ell_{23}&=m_1\Delta_2(r_{31}^{-3}-r_{12}^{-3}),\\
 \dot\ell_{31}&=m_2\Delta_2(r_{12}^{-3}-r_{23}^{-3}).
\end{aligned}
\end{equation}
Their derivatives have signs $(-,+,-)$ at launch in the strict fundamental
interval. Since every $\ell_{ij}$ vanishes at a brake, each right-hand
side must integrate to zero before a second brake. Its initial strict
sign must therefore be reversed somewhere. Before the first syzygy, this
requires reversals of all three initial side comparisons, at possibly
different times. It does not prove that those reversals cannot occur.

Integrating Lagrange--Jacobi between two brakes also gives
\[
 \int_0^\tau U\,dt=2U_0\tau,\qquad
 \int_0^\tau K\,dt=U_0\tau.
\]
Integrating the labelled accelerations against the fixed launch positions
also gives
\[
 \int_0^\tau\sum_{i<j}m_im_j
 {d_{ij}(0)\cdot d_{ij}(t)\over r_{ij}(t)^3}\,dt=0.
\]
The integrand is initially positive; its later sign is not controlled.

\subsection{The torque ratio and its threshold}

On the arc from launch to the first loss of positive area or strict ordering
$r_{12}>r_{23}>r_{31}$, put
$R=r_{12}$, $x=r_{23}/R$, $y=r_{31}/R$, and define
\[
 F=\Delta_2(r_{31}^{-3}-R^{-3})>0,\qquad
 G=\Delta_2(r_{23}^{-3}-R^{-3})>0,\qquad
 k={G\over F}={y^3(1-x^3)\over x^3(1-y^3)}.
\]
Integrating the torques from release gives
\begin{equation}\label{eq:history}
 \eta={\int_0^tG\,d\tau\over\int_0^tF\,d\tau}
 ={m_1\over m_2}{-\ell_{31}\over\ell_{23}},\qquad
 \dot\eta={F\over\int_0^tF\,d\tau}(k-\eta).
\end{equation}
Thus $\eta$ is a weighted average of the preceding values of $k$.
Strict decrease of $k$ would imply $k(t)<\eta(t)<k(0)$; establishing that
decrease along the tied trajectories is the remaining problem.

Use shape time $ds=dt/R^{3/2}$ and write
$W=\ell_{23}/\sqrt R$, $\delta=\Delta_2/R^2$. Eliminating the two
shape velocities gives the rational identity
\begin{equation}\label{eq:threshold}
 (\log k)_s={W\over\delta}\mathcal A(u,x,y)(\eta-h(u,x,y)),
 \qquad \mathcal A>0,\quad h>k.
\end{equation}
For an explicit definition, put $m=m_1$, $n=m_2$, and
\[
\begin{aligned}
 \mathsf B={}&n(-x^5-x^3y^2+x^3+x^2+2y^5-y^2-1)\\
 &+x^3y^2(x^2-y^2+1)+y^5(-x^2+y^2+1)-2y^2,\\
 \mathsf N={}&m(2x^5-x^2y^3-x^2-y^5+y^3+y^2-1)\\
 &+x^5(x^2-y^2+1)+x^2y^3(-x^2+y^2+1)-2x^2,\\
 h={}&{y^2\mathsf N\over x^2\mathsf B},\qquad
 \mathcal A={-3\mathsf B\over
 2my^2(x-1)(y-1)(x^2+x+1)(y^2+y+1)}.
\end{aligned}
\]
Here $h$ denotes the torque threshold. The ordered-triangle substitution
$x=1-ab/2$, $y=1-a+ab/2$, $u=(\sqrt2-1)v$, with $0<a,b,v<1$,
maps the domain to a cube. After clearing denominators with known sign,
exact tensor-Bernstein coefficients prove $\mathcal A>0$ and $h-k>0$.
Thus~\eqref{eq:threshold} reduces the desired monotonicity to
$\eta<h$. Section~\ref{sec:iso} proves this inequality on the first
arc near the equal-leg endpoint.

\subsection{A corrected conditional sign-change argument}

Define the gap, lag, and threshold distance
\[
 d=h-k>0,\qquad e=\eta-k,\qquad w=h-\eta=d-e,
\]
and the positive coefficients
\[
 \lambda={m_1\delta(y^{-3}-1)\over W},\qquad
 c={kW\mathcal A\over\delta}.
\]
Equations~\eqref{eq:history}--\eqref{eq:threshold} give
\begin{equation}\label{eq:lag}
 e_s+(\lambda+c)e=cd,\qquad
 w_s+(\lambda+c)w=f,\qquad f=d_s+\lambda d.
\end{equation}
For $s$ strictly before syzygy, the launch condition has the improper
integral representation
\[
 e(s)=\lim_{\varepsilon\downarrow0}\int_\varepsilon^s c(\tau)d(\tau)
 \exp\!\left[-\int_\tau^s(\lambda+c)\,d\xi\right]d\tau.
\]
At launch $W$ has a simple zero, $\lambda\sim1/s$, and $e(0)=0$ after
analytic extension, so the homogeneous launch contribution tends to zero.
For a fixed interior reference point $s_0>0$, define the integrating factor
\[
 M(s)=\exp\!\int_{s_0}^s(\lambda+c)\,d\xi,
 \qquad (Mw)_s=Mf.
\]

At the first transverse ordered syzygy $s=S$, collinear momentum
constraints force $w(S)=0$, while $\delta\to0$ and
$c=kW\mathcal A/\delta$ can diverge. The endpoint equality therefore
does not imply $Mw\to0$, as assumed in an earlier draft.

\begin{lemma}[Sufficient sign-change criterion with terminal control]
Assume $w>0$ near launch, $f$ is positive up to one interior sign change
and negative thereafter, and
\begin{equation}\label{eq:terminal-weighted}
 \liminf_{s\uparrow S} M(s)w(s)\ge0.
\end{equation}
Then $w>0$ throughout $0<s<S$.
\end{lemma}
\begin{proof}
Start at any sufficiently small positive time where $w>0$. Since $M>0$,
the product $Mw$ increases until the switch and then strictly decreases.
If it reached zero before $S$, its terminal lower limit would be negative,
contrary to~\eqref{eq:terminal-weighted}.
\end{proof}

The missing hypothesis cannot be omitted from this scalar argument.
On $0<s<1$ take
\[
 a={1\over s}+{1\over1-s},\qquad
 w=(1-s)(1+s-3s^2),\qquad
 f=(1-s)\left({1\over s}+2-9s\right).
\]
Then $w_s+aw=f$, $w(0)=1$, and $w(1)=0$. The forcing changes from positive
to negative exactly once, at $(1+\sqrt{10})/9$, yet $w$ becomes negative
after $(1+\sqrt{13})/6$. Up to a positive constant,
$M=s/(1-s)$ and $Mw\to-1$.
This scalar counterexample shows why terminal control is necessary.
It makes no claim about a Newtonian trajectory.

\subsection{Amplitude bounds and the terminal face}

The two centrifugal contributions to the side-gap accelerations, after
multiplication by $R^2$, are
\[
 Z\bigl((1-\eta)^2/m^2-x^{-3}\bigr),\qquad
 Z\bigl(x^{-3}-n^2\eta^2/(m^2y^3)\bigr),\qquad Z=W^2.
\]
Thus shape and the two history variables $(Z,\eta)$ determine both terms.
At a contact $\eta=h$, put $X_c=\delta x_s/W$,
$Y_c=\delta y_s/W$,
$S_c=h_xX_c+h_yY_c$, and
$P_c=m\delta^2(y^{-3}-1)(k-h)<0$. Direct differentiation gives
\[
 W\delta(\eta-h)_s=P_c-ZS_c.
\]
Contacts with $S_c\ge0$ point into $\eta<h$. On $S_c<0$, the
sufficient amplitude bound is $Z<Z_J:=P_c/S_c$.

At an ordered syzygy with body 3 between bodies 1 and 2, write
$y=q$, $x=1-q$. The launch angular-momentum signs require
$0<q<n/(m+n)$, and the collinear momentum constraints give
\[
 h=\eta_{\rm col}={q(m+1-q)\over(1-q)(n+q)}.
\]
Canceling the common area-squared factor in $P_c/S_c$ defines the
terminal threshold $Z_J^{\rm syz}$. Exact Bernstein certificates prove
$1<Z_J^{\rm syz}<Z_2$, where
$Z_2=R^2(\mathfrak g_{23}-\mathfrak g_{31})/
[n^2\eta_{\rm col}^2/(m^2q^3)-(1-q)^{-3}]$ is the amplitude at which
the second-gap acceleration vanishes, and
$\mathfrak g_{ij}=\ddot r_{ij}-\ell_{ij}^2/r_{ij}^3$.
The denominator is positive on this face. At the equal-leg endpoint the
first-syzygy certificate gives $Z_*<1$; its tied tangent satisfies
$\partial_v[q(m+n)/n]>18$. These bounds put nearby $v<1$ members
strictly inside the torque-compatible face with $Z<1$.
For the full family, neither interior noncontact nor the terminal
amplitude bound is proved.

Energy and angular momentum alone do not imply the required bounds.
For masses $(m,n,1)=(4/5,3/5,1)$, consider an ordered collinear state with
\[
 (r_{12},r_{23},r_{31})=(1,697/700,3/700),\qquad
 \eta=419/32759,\qquad
 \ell_{23}^2={476958362267\over2867548600}.
\]
Choose zero radial velocities and the transverse velocities fixed by these
angular momenta and $P=L=0$. Exact arithmetic gives $H=-U_0$, pair angular-momentum
signs $(-,+,-)$, and first-crossing area orientation
$\dot\Delta_2/\ell_{23}=-292607/131036<0$, but
\[
 (r_{23}-r_{31})''=-{96354167469624287\over1827657180810}<0.
\]
This state is not asserted reachable from the tied initial brake. It shows
that energy, collinear geometry, crossing direction, and pair angular-momentum signs do
not suffice to control the second distance gap. The main analytic target
is therefore a bound on \emph{reachable histories}.
